\hypertarget{python-3.7-dataclasses-context-variables-and-dict-ordering-guarantees}{%
\chapter{Python 3.7: Dataclasses, Context Variables, and Dict Ordering
Guarantees}\label{python-3.7-dataclasses-context-variables-and-dict-ordering-guarantees}}

\hypertarget{pep-557-dataclasses-under-the-hood}{%
\subsection{13.1 PEP 557: Dataclasses under the
hood}\label{pep-557-dataclasses-under-the-hood}}

Introduced in Python 3.7, the \passthrough{\lstinline!@dataclass!}
decorator automates boilerplate method generation. * \textbf{Code
Generation}: Rather than using dynamic wrappers at runtime,
\passthrough{\lstinline!@dataclass!} is a class decorator that runs at
import time. It reads the class's
\passthrough{\lstinline!\_\_annotations\_\_!} dictionary and dynamically
generates the source code for methods like
\passthrough{\lstinline!\_\_init\_\_!},
\passthrough{\lstinline!\_\_repr\_\_!}, and
\passthrough{\lstinline!\_\_eq\_\_!}. It then compiles and attaches
these methods to the class dict, ensuring runtime performance is
identical to manually written methods.

\hypertarget{properties-cached-properties}{%
\subsection{13.2 Properties \& Cached
Properties}\label{properties-cached-properties}}

\begin{itemize}
\tightlist
\item
  \textbf{The Property Descriptor}: The built-in
  \passthrough{\lstinline!@property!} decorator is a data descriptor. It
  intercepts access to attributes and runs custom getter, setter, and
  deleter methods:
\end{itemize}

\begin{lstlisting}[language=Python]
class Circle:
    def __init__(self, radius):
        self._radius = radius

    @property
    def radius(self):
        return self._radius

    @radius.setter
    def radius(self, value):
        if value < 0:
            raise ValueError("Radius cannot be negative")
        self._radius = value
\end{lstlisting}

\begin{itemize}
\tightlist
\item
  \textbf{Cached Properties}: For expensive computations,
  \passthrough{\lstinline!functools.cached\_property!} caches results.
  It operates as a non-data descriptor:

  \begin{enumerate}
  \def\labelenumi{\arabic{enumi}.}
  \tightlist
  \item
    The first access computes the value and writes it directly to the
    instance's dictionary \passthrough{\lstinline!\_\_dict\_\_!}.
  \item
    Subsequent lookups find the value in
    \passthrough{\lstinline!\_\_dict\_\_!} directly, bypassing the
    descriptor's getter logic.
  \end{enumerate}
\end{itemize}

\hypertarget{instance-slots-optimization}{%
\subsection{13.3 Instance Slots
Optimization}\label{instance-slots-optimization}}

By default, every instance allocates a dictionary
\passthrough{\lstinline!\_\_dict\_\_!} to store attributes. This allows
dynamic attribute addition but consumes significant memory. *
\textbf{Slots (\passthrough{\lstinline!\_\_slots\_\_!})}: Defining
\passthrough{\lstinline!\_\_slots\_\_!} inside a class tells CPython to
use a fixed-size array instead of a dictionary to store attributes. *
\textbf{CPython Layout}: The class MRO creates descriptor references
mapping each slot to a static array offset in C. This: 1. Eliminates the
memory overhead of the instance dictionary
\passthrough{\lstinline!\_\_dict\_\_!} and
\passthrough{\lstinline!\_\_weakref\_\_!}. 2. Speeds up attribute access
by replacing dictionary key lookups with fast array indexing. 3.
Prevents users from dynamically adding arbitrary new attributes to the
instance.

\begin{lstlisting}[language=Python]
import sys

class DictClass:
    def __init__(self, x, y):
        self.x = x
        self.y = y

class SlotsClass:
    __slots__ = ('x', 'y')
    def __init__(self, x, y):
        self.x = x
        self.y = y

d = DictClass(1, 2)
s = SlotsClass(1, 2)

print("DictClass size:", sys.getsizeof(d) + sys.getsizeof(d.__dict__))
# SlotsClass has no __dict__, consuming significantly less memory
print("SlotsClass size:", sys.getsizeof(s))
\end{lstlisting}

\begin{center}\rule{0.5\linewidth}{0.5pt}\end{center}
